% THE Q-MODEL — From Q to Cosmology
% Version 1.0 (revised) — White Paper (Zenodo release)
% Build: pdflatex TQM_v1_0.tex   (twice for references, if any)

\documentclass[11pt,a4paper]{article}
\usepackage[utf8]{inputenc}
\usepackage[T1]{fontenc}
\usepackage{amsmath,amssymb,amsthm}
\usepackage[margin=2.5cm]{geometry}
\usepackage{booktabs}
\usepackage{lmodern}
\usepackage[hidelinks]{hyperref}
\usepackage{xcolor}

\definecolor{caveat}{HTML}{7A2E2E}

\title{\bfseries THE Q-MODEL — From Q to Cosmology\\[2mm]
\large A Theory of Structure, Complexity and Random Actualization}
\author{Fabrice Wieser}
\date{Version 1.0 (revised) \quad·\quad 15 August 2026}

\begin{document}

\maketitle

\begin{center}
\fcolorbox{caveat}{white}{%
\begin{minipage}{0.92\textwidth}
\color{caveat}
\textbf{Publication status: READY\_FOR\_WHITEPAPER — NOT\_READY\_FOR\_JOURNAL.}\\[2mm]
This document is released as a \emph{white paper} (research-program description) for
archival on Zenodo. It is \emph{not} submitted for peer-reviewed journal publication. The
central derivation claim (Einstein recovery from Q-events) remains \emph{logical, not
mathematical}: the metric and the BDG action are imported (proven but not TQM-derived),
$G=\ell^2c^3/\hbar$ is dimensional analysis, and no unique sharp prediction yet
discriminates TQM from SM$+\Lambda$CDM. See
\texttt{Docs/Audits/PublicationReadiness\_Final.md}.
\end{minipage}}
\end{center}

\vspace{2mm}

\begin{abstract}
\noindent
THE Q-MODEL (TQM) is a theory of structure and content that compresses observable physics
to four primitives — the individuation principle $Q$, Random Actualization, the scale
triad $(\ell,\tau,\hbar)$, and a single continuous nonlinearity parameter $M^2$ — and
holds that the \emph{form} of every structure is \textbf{derivable} from these primitives
while the \emph{content} (specific masses, couplings, multiplicities, the Koide angle) is
\textbf{realized} by contingent draw. The derivation program yields $U(1)$ as a theorem
($\mathrm{Aut}(S^1)=U(1)$, $\pi_1(S^1)=\mathbb{Z}$; 0.95), spatial dimensionality
$d=3{+}1$ from the intersection of physical viability windows ($M^2\approx5$; 0.85), the
multiplicity lower bound $N\ge3$ from CP violation, the log-normal abundance law from the
central limit theorem, and phase-gradient gravity whose leading order is the Einstein
equations. The dynamical layer is the graph Laplacian $L_Q$, whose tight-binding identity
$H=tL_Q$ and reversible dynamics yield the Schr\"odinger equation
$i\partial_t\psi=L_Q\psi$ and a Hilbert space of eigenmodes. The non-abelian gauge factors
and the Koide relation are classified by a minimal three-category taxonomy
(DERIVED / REAL-UNDERIVED / DRAWN) with program consistency 0.95 and zero category
conflicts. The theory makes quantitative, falsifiable predictions — the zero-parameter
radial acceleration $g_\dagger=cH_0/(2\pi)$, a specific $w(z)$, a time-varying
$\Lambda(t)$, and the log-normal abundance form — and one such prediction (neutrino-Koide)
was already falsified. Five no-go theorems (T-08 through T-12) are conditional statements
under the no-new-primitives constraint.
\end{abstract}

\noindent\textbf{Keywords:} THE Q-MODEL; TQM; emergence; quantum foundations; causal set
theory; graph Laplacian; Schr\"odinger equation; gauge structure; radial acceleration
relation; dark energy; structure/content split; falsifiability.

\tableofcontents

\section{Introduction}
\label{sec:intro}

The Standard Model describes the \emph{content} of particle physics but leaves its
\emph{values} unexplained. THE Q-MODEL (TQM) addresses this with a single organizational
claim: \textbf{structure is derivable; content is realized.}

Every result is classified by a minimal three-category taxonomy:
\begin{itemize}
\item \textbf{DERIVED} — computable from the primitives by theorem;
\item \textbf{REAL-UNDERIVED} — real, precise, predictive structure whose origin is not
      computable (emergent, or structured);
\item \textbf{DRAWN} — a coincidental draw of the abundance law.
\end{itemize}
The taxonomy classifies \emph{underivability}; REAL-UNDERIVED and DRAWN are \emph{not}
claimed as derivations.

\section{Formal Primitive Definitions}
\label{sec:primitives}

\subsection{Ontology layer}
\begin{center}
\begin{tabular}{ll}
\toprule
Primitive & Role \\
\midrule
$Q$ & principle of individuation (ontology) \\
Random Actualization & genuine ontological chance — \textbf{assumption A-03}, not derived \\
$(\ell,\tau,\hbar)$ & spacetime scale, clock, action (unit conventions) \\
$M^2$ & nonlinearity regime (single contingent continuous parameter) \\
\bottomrule
\end{tabular}
\end{center}

\subsection{Dynamical layer (quantum postulates)}
\begin{enumerate}
\item \textbf{$Q$ exists.} Topological charge quanta have position $x_i$ and phase
      $\theta_i$, and interact pairwise with coupling $J_{ij}=f(|x_i-x_j|)$.
\item \textbf{Reversible dynamics.} The state norm $\|\psi\|^2$ is conserved
      ($\Leftrightarrow$ unitarity).
\item \textbf{Born rule.} $P=|\langle\phi|\psi\rangle|^2$, uniquely selected by additivity
      (Gleason's theorem).
\item \textbf{Measurement.} A measurement yields one definite outcome (collapse axiom).
\end{enumerate}

\section{Dynamical System Summary}
\label{sec:dynamics}

The dynamical content of TQM is the graph Laplacian $L_Q=D-A$ (adjacency $A$, degree $D$):
real symmetric, positive semi-definite, zero row sums. Its eigenvectors form the
\emph{Hilbert space} of the theory. The tight-binding identity (TQM-142) is
$(L_Q)_{ij}=2\delta_{ij}-\delta_{i,j+1}-\delta_{i,j-1}$, giving $H=tL_Q$.

\textbf{Schr\"odinger from reversibility (TQM-149--151).} For linear evolution
$\partial_t\psi=M\psi$, norm conservation forces $M^\dagger=-M$; with the complex
structure $J$ ($J^2=-I$), $M=J\otimes L_Q$ gives
\begin{equation}
i\partial_t\psi = L_Q\,\psi,\qquad \psi(t)=e^{-iL_Qt}\psi(0).
\end{equation}

\textbf{Continuum limit (verified).} The discrete chain converges to the flat Laplacian
with exact spectrum $\lambda_k=(1/\Delta x^2)\bigl[2-2\cos(\pi k/(N{+}1))\bigr]\to(\pi k)^2$;
the BDG operator converges to the d'Alembertian $\square$ with $O(h^2)$ error. The weighted
Laplacian $L_W=D_K-K$ approximates the Laplace--Beltrami operator $\Delta_g$ and defines a
unitary curved-space Schr\"odinger equation. (See \texttt{TQM.Tests/ResearchXC}.)

\section{Structure/Content Split}
\label{sec:split}

\textbf{Form} (topology, symmetry, distribution shape, lower bounds) is \textbf{derived};
\textbf{content} (drawn values, specific angles, upper bounds) is \textbf{contingent}. The
split is \emph{falsifiable} (§\ref{sec:predictions}): one prediction (neutrino-Koide) was
falsified. Two composite objects are admitted: internal $N=3$ =
DERIVED (lower bound) $\cap$ DRAWN (upper bound); $SU(3)$ whole =
REAL-UNDERIVED (structure) + DRAWN (count 3).

\section{Derivation Hierarchy}
\label{sec:hierarchy}

\begin{center}
\begin{tabular}{ll}
\toprule
Chain & Result \\
\midrule
Complexity maximization & $M^2\approx5\Rightarrow d=3{+}1$ (0.85) \\
Oscillation $\to$ phase $\to$ $S^1$ vacuum & $U(1)$ (0.95) \\
Binary defects $\to$ $Z_2$ winding & $SU(2)$ structure (0.70) \\
Tri-defect moduli & $SU(3)$ structure (0.10); count DRAWN \\
\bottomrule
\end{tabular}
\end{center}

\section{Complexity Functional}
\label{sec:complexity}

The complexity argument is a \textbf{weighted six-component decomposition} of observer
viability (\texttt{ComplexityOptimumAnalyzer.cs}), not a single scalar variational
functional. $d=3{+}1$ is the unique value satisfying Bertrand's theorem, Gauss's law,
knot stability, Huygens, and 2 GR polarizations. $M^2\approx5$ is the
observer-viability intersection. \emph{Honesty note:} a window-intersection argument, not
a variational theorem (T-02, 0.85).

\section{Gauge Sector}
\label{sec:gauge}

$U(1)$ is DERIVED ($\mathrm{Aut}(S^1)=U(1)$, $\pi_1(S^1)=\mathbb{Z}$; T-01). $SU(2)$ and
$SU(3)$ structure are REAL-UNDERIVED (emergent); the color count $n=3$ is DRAWN. TQM
derives the gauge-group \emph{structure}, \emph{not} the gauge \emph{dynamics} (Maxwell /
Yang--Mills actions are borrowed, A-07).

\section{Flavor Sector}
\label{sec:flavor}

The Yukawa operator is the overlap operator $Y_{ij}=\langle
\mathrm{arch}_i|\mathrm{amplitude}|\mathrm{arch}_j\rangle$ (form derived, spectrum
underived). The Koide relation $Q=2/3$ at $\theta=45^\circ$ is \textbf{REAL-UNDERIVED,
structured, closed (underivable)} (T-08, 0.70). Neutrino-Koide is \textbf{FALSIFIED}
(T-11, 0.90): $Q_{\max}=0.585$ (normal) / $0.500$ (inverted) $<2/3$.

\section{Gravity and Emergent GR}
\label{sec:gravity}

$G=\ell^2c^3/\hbar$ is a \emph{dimensional} result, not a derivation of gravity's
dynamics. The phase-gravity chain (oscillation $\to$ phase $\to$ causal density $\to$
metric $\to$ curvature $\to$ Einstein) yields
$G_{\mu\nu}=8\pi G_{\rm eff}T_{\mu\nu}+O(\ell_P^2 R^2)$ at leading order, with two
modifications beyond GR: $\Lambda(t)=\alpha/\sqrt{V(t)}$ and singularity resolution at
$r\sim\ell_P$.

\emph{Honesty note:} the chain is \textbf{logical, not a new dynamical derivation} —
"same equations, same predictions; no falsifiable difference" in the weak field. The value
is \emph{ontological}; the physical content beyond GR is confined to $\Lambda(t)$ and
singularity resolution. The RAR $g_\dagger=cH_0/(2\pi)$ is derived with \emph{zero free
parameters} (the $2\pi$ is a numerical accident).

\section{Cosmology}
\label{sec:cosmology}

$\Lambda\sim1/\sqrt{N}$ is a postdiction; $w(z)=-1+0.015(1+z)^{3/2}$; clusters are
reproduced; the CMB is an \emph{accepted partial computational layer} (full $C_\ell$ solver
deferred, not a theory gap).

\section{Quantitative Predictions}
\label{sec:predictions}

\begin{center}
\begin{tabular}{lll}
\toprule
Prediction & Type & Status \\
\midrule
RAR $g_\dagger=cH_0/(2\pi)\approx1.05\times10^{-10}$ & zero-parameter & matches $a_0$ \\
$w(z)=-1+0.015(1+z)^{3/2}$ & cosmological & Euclid $>3\sigma$ by $\sim$2030 \\
$\Lambda(t)=\alpha/\sqrt{V(t)}$ & dark energy & Euclid / Roman \\
log-normal abundance law & distribution form & testable \\
$N\ge3$ (CP lower bound) & a priori & confirmed \\
neutrino-Koide $Q=2/3$ & falsifiable & \textbf{falsified} \\
\bottomrule
\end{tabular}
\end{center}

\section{Classification System, No-Gos, Open Questions}
\label{sec:closure}

The minimal taxonomy classifies fourteen results (consistency 0.95, zero category
conflicts, two composites). Five \textbf{conditional no-go theorems} (T-08--T-12) bound
what the theory cannot do under the no-new-primitives constraint. Every in-scope question
is \textbf{dispositioned} (derived, underivable, underdetermined, contingent, or accepted
as a computational layer); the internal-3 node is dispositioned
\emph{unresolved-contingent} (gauge-count no-go T-09 provisional at 0.10).

\section{Scope and Limitations}
\label{sec:scope}

TQM derives the \emph{form} of structure and \emph{classifies} content. It does \emph{not}
derive gauge dynamics (A-07), nor contingent values (masses, couplings, multiplicities,
Koide angle). Confidence numbers are uncalibrated ordinal ranks. The phase-gravity chain is
ontological in the weak field. The unified action is a roadmap, not a result.

\section{Conclusion}
\label{sec:conclusion}

THE Q-MODEL is \emph{structurally complete} in the precise sense that every in-scope
question is \textbf{dispositioned}. The single residual of genuine scientific tension is
the \textbf{internal-3 node} (why internal multiplicity/count saturates at 3),
dispositioned unresolved-contingent with a provisional gauge-count no-go (T-09, 0.10).
\textbf{Structure is derivable. Content is realized.}

\end{document}
